\uuid{w7Oi}
\titre{Équation droite/plan}
\auteur{bodin ; hanani ; mzari}
\organisation{exo7}
\licence{CC BY-NC-SA 4.0}
\datecreate{2020-01-20}
\niveau{L1}
\module{Géométrie}
\chapitre{Géométrie du plan}
\sousChapitre{Distance, norme, produit scalaire}
\difficulte{2}
\theme{équation droite, plan}
\exo7qcmid{lille-1.geometrie-plane.016}

\contenu{
\begin{question}
\qclassification{140.01,140.02}
\qtags{motcle=équation droite/plan}
Soit $D$ la droite définie par le paramétrage :
$$\left\{\begin{array}{ccl}x&=&1+t\\ y&=&2-t,\quad t\in \Rr.
\end{array}\right.$$
Quelles sont les bonnes réponses ?
\begin{answers}
    \bad{Le point $A(1,1)$ appartient à $D$.}
    \bad{$\vec{u}=(1,-1)$ est un vecteur normal à $D$.}
    \good{Une équation cartésienne de $D$ est : $x+y-3=0$.}
    \bad{$\vec{u}(1,1)$ est un vecteur directeur de $D$.}
\end{answers}
\begin{explanations}
Le vecteur $\vec{u}=(1,-1)$ est un vecteur directeur de $D$. On élimine $t$ en additionnant les deux équations. Ceci donne $x+y=3$ qui est une équation cartésienne de $D$.
\end{explanations}
\end{question}
}
